\chapter{Μηχανικές Μετρήσεις της Διαδρομής H48}
\label{app:h48eng}

Το παράρτημα αυτό συγκεντρώνει τις μηχανικές (\eng{engineering}) μετρήσεις της διαδρομής H48, οι οποίες αφορούν το κόστος φόρτωσης, ελέγχου και επαναχρησιμοποίησης του πίνακα αποκοπής και των σχετικών διεργασιών, και όχι τη δυσκολία της ίδιας της αναζήτησης. Οι μετρήσεις τεκμηριώνουν τις σχεδιαστικές επιλογές της υλοποίησης (έμπιστος πίνακας, λειτουργία δέσμης, παραμένουσα διεργασία, διεπαφές), αλλά δεν στηρίζουν κανέναν ισχυρισμό βελτιστότητας πέρα από όσα παρουσιάζονται στο Κεφάλαιο~\ref{chap:experiments}. Τα συμπεράσματά τους συνοψίζονται εκεί σε μία παράγραφο· εδώ δίνεται η πλήρης καταγραφή.

Τα χρονικά όρια των μετρήσεων του παραρτήματος είναι τα εξής: 180~s για την έμπιστη πιστοποίηση με προφόρτωση, 300~s για την έμπιστη πιστοποίηση χωρίς προφόρτωση, 240~s για την πιστοποίηση με παραμένουσα διεργασία, 120~s για τη σύγκριση ελεγχόμενης/έμπιστης λειτουργίας, 180~s για τη λειτουργία δέσμης, 30~s για την παραμένουσα διεργασία και 60~s για τις διεπαφές γραμμής εντολών, ροής και προγραμματιστικής χρήσης. Τα κυριότερα δίκαια συγκριτικά τεκμήρια είναι τα \path{results/processed/h48_trusted_table_speedup_seed_2026_thesis_h48h7_fair.json}, \path{results/processed/h48_batch_overhead_seed_2026_thesis_trusted_fair.json} και \path{results/processed/h48_resident_oracle_seed_2026_thesis_h48h7_trusted_fair.json}.

\section{Δίκαιη λογιστική των επιταχύνσεων}
\label{sec:experiments-h48-fair}

Η διαδρομή H48 υποστηρίζει δύο τρόπους χρήσης του παραγόμενου πίνακα: την ελεγχόμενη λειτουργία (\eng{checked}), όπου κάθε κλήση επαληθεύει πλήρως τον πίνακα με \code{nissy_checkdata}, και την έμπιστη λειτουργία (\eng{trusted}), όπου ο πίνακας γίνεται δεκτός με βάση τα μεταδεδομένα παραγωγής του. Η έμπιστη λειτουργία στηρίζεται στα μεταδεδομένα του πίνακα \code{h48h7}· το ζήτημα της μη ντετερμινιστικής παραγωγής του πίνακα αυτού και η τεκμηριωμένη απόφαση διατήρησης των σχετικών τεκμηρίων συζητούνται στο Κεφάλαιο~\ref{chap:validation}, μαζί με την ανεξάρτητη διασταύρωση όλων των βαθιών ακριβών γραμμών.

Απαιτείται μια σημείωση μεθοδολογίας για τη λογιστική των μετρήσεων. Σε πρώτη καταγραφή χωρίς κοινή προθέρμανση, οι αποθηκευμένοι πολλαπλασιαστές ήταν 189.894x για τη σύγκριση ελεγχόμενης/έμπιστης λειτουργίας, 11.56x για τη δέσμη με ελεγχόμενο πίνακα, 38.198x για τη δέσμη με έμπιστο πίνακα και 40.16x για την παραμένουσα διεργασία. Οι τιμές αυτές χρέωναν το εφάπαξ ψυχρό κόστος εισαγωγής σελίδων (\eng{cold page-in}) του πίνακα των 3\,793\,842\,344 bytes εξ ολοκλήρου στο αργό σκέλος κάθε σύγκρισης. Τα αναθεωρημένα δίκαια τεκμήρια εκτελούν προθέρμανση πριν από κάθε χρονομετρημένο σκέλος και καταγράφουν σταθερή σειρά σκελών, ώστε κανένα σκέλος να μην πληρώνει το εφάπαξ κόστος μέσα στο μετρούμενο παράθυρο. Όλες οι τιμές που ακολουθούν προέρχονται από τα δίκαια τεκμήρια.

Στη δίκαιη σύγκριση ελεγχόμενης και έμπιστης λειτουργίας (Πίνακας~\ref{tab:h48-trusted-fair}), οι τέσσερις ελεγχόμενες και οι τέσσερις έμπιστες κλήσεις είναι όλες \code{exact} και επαληθευμένες: 42.058193~s για το ελεγχόμενο σκέλος έναντι 0.192174~s για το έμπιστο, δηλαδή επιτάχυνση 218.855x συνολικά και 232.606x σε σταθερή κατάσταση, εξαιρώντας την πρώτη χρονομετρημένη κλήση κάθε σκέλους. Η αναλογία είναι γνήσια και όχι τέχνημα κρυφής μνήμης, επειδή το ελεγχόμενο σκέλος επανασαρώνει ολόκληρο τον πίνακα \code{h48h7} σε κάθε κλήση. Η μέτρηση δείχνει ότι η καθυστέρηση βρίσκεται κυρίως στη σάρωση και επαλήθευση του μεγάλου πίνακα ανά διεργασία· δεν πρέπει να διαβαστεί ως επιτάχυνση αναζήτησης χειρότερης περίπτωσης για όλο τον χώρο καταστάσεων.

\begin{table}[htbp]\centering
{\small
\begin{tabular}{lrrr}
\hline
Mode & Total (s) & Steady-state (s) & Exact \\
\hline
Verified each call & 42.058193 & 31.739583 & 4 \\
Trusted generated table & 0.192174 & 0.136452 & 4 \\
Speedup & 218.855x & 232.606x & -- \\
\hline
\end{tabular}
}

\caption[Δίκαιη σύγκριση ελεγχόμενης και έμπιστης χρήσης του πίνακα h48h7]{Δίκαιη σύγκριση ελεγχόμενης και έμπιστης χρήσης του πίνακα \code{h48h7}, με κοινή προθέρμανση ώστε κανένα σκέλος να μη χρεώνεται το εφάπαξ ψυχρό κόστος φόρτωσης. Πηγή: \path{results/processed/h48_trusted_table_speedup_seed_2026_thesis_h48h7_fair.json}.}
\label{tab:h48-trusted-fair}
\end{table}

Η αντίστοιχη δίκαιη μέτρηση της λειτουργίας δέσμης με έμπιστο πίνακα (Πίνακας~\ref{tab:h48-batch-fair}) δείχνει ότι οι 12 επαναλήψεις της ίδιας ρηχής περίπτωσης παραμένουν \code{exact} και επαληθευμένες, με 0.67774~s σε ξεχωριστές έμπιστες κλήσεις έναντι 0.181076~s σε μία δέσμη, δηλαδή ρυθμαπόδοση (\eng{throughput}) 3.743x συνολικά και 3.112x σε σταθερή κατάσταση. Αφού αφαιρεθούν ο πλήρης έλεγχος του πίνακα και το εφάπαξ κόστος φόρτωσης, το πραγματικό κόστος διεργασίας και δέσμης για μικρές καταστάσεις παραμένει μετρήσιμο, αλλά είναι πολύ μικρότερο από τον αρχικά καταγεγραμμένο πολλαπλασιαστή.

\begin{table}[htbp]\centering
{\small
\begin{tabular}{lrrr}
\hline
Mode & Total (s) & Steady-state (s) & Exact \\
\hline
Single process per solve & 0.67774 & 0.516481 & 12 \\
Batch, table loaded once & 0.181076 & -- & 12 \\
Throughput speedup & 3.743x & 3.112x & -- \\
\hline
\end{tabular}
}

\caption[Δίκαιη μέτρηση του κόστους δέσμης με έμπιστο πίνακα]{Δίκαιη μέτρηση του κόστους δέσμης με έμπιστο πίνακα \code{h48h7} (κοινή προθέρμανση, 12 επαναλήψεις): ξεχωριστές διεργασίες έναντι μίας δέσμης με τον πίνακα φορτωμένο μία φορά. Πηγή: \path{results/processed/h48_batch_overhead_seed_2026_thesis_trusted_fair.json}.}
\label{tab:h48-batch-fair}
\end{table}

Τέλος, η δίκαιη σύγκριση ξεχωριστών κλήσεων και παραμένουσας διεργασίας (\eng{resident process}) στον Πίνακα~\ref{tab:h48-resident-fair} δίνει 0.325648~s για τις 8 ξεχωριστές κλήσεις έναντι 0.085528~s για τις 8 κλήσεις της παραμένουσας διεργασίας, δηλαδή 3.808x συνολικά και 6.55x σε σταθερή κατάσταση, εξαιρώντας την πρώτη κλήση κάθε σκέλους και την εκκίνηση της συνεδρίας. Η μέτρηση αφορά την επαναλαμβανόμενη χρήση και την αποφυγή της επανεκκίνησης διεργασίας και της επανεγκατάστασης του πίνακα για κάθε ρηχή κατάσταση.

\begin{table}[htbp]\centering
{\small
\begin{tabular}{lrrr}
\hline
Mode & Total (s) & Steady-state (s) & Exact \\
\hline
Separate native process & 0.325648 & 0.274602 & 8 \\
Resident native process & 0.085528 & 0.041924 & 8 \\
Speedup & 3.808x & 6.55x & -- \\
\hline
\end{tabular}
}

\caption[Δίκαιη σύγκριση ξεχωριστών και παραμένουσας εγγενούς διεργασίας]{Δίκαιη σύγκριση ξεχωριστών εγγενών διεργασιών και παραμένουσας εγγενούς διεργασίας για επαναλαμβανόμενες κλήσεις με έμπιστο πίνακα \code{h48h7} (κοινή προθέρμανση, 8 επαναλήψεις). Πηγή: \path{results/processed/h48_resident_oracle_seed_2026_thesis_h48h7_trusted_fair.json}.}
\label{tab:h48-resident-fair}
\end{table}

\section{Παραλλαγές πιστοποίησης του δύσκολου σώματος}
\label{sec:h48eng-cert-variants}

Η πιστοποίηση του δύσκολου σώματος σε ελεγχόμενη λειτουργία παρουσιάζεται στο Κεφάλαιο~\ref{chap:experiments} (Πίνακας~\ref{tab:h48-cert}). Η ίδια πιστοποίηση εκτελέστηκε και σε έμπιστη λειτουργία με προφόρτωση του πίνακα (Πίνακας~\ref{tab:h48-cert-preload}). Το αποτέλεσμα παραμένει 4/4 \code{exact} και επαληθευμένο, αλλά η γραμμή του superflip χρειάστηκε 158.440286~s. Πρόκειται για σημαντικό αρνητικό εύρημα για την ερμηνεία των επιταχύνσεων: η έμπιστη και προφορτωμένη διαδρομή αποφεύγει τον επαναλαμβανόμενο έλεγχο του πίνακα, αλλά δεν καθιστά τη δυσκολότερη αναζήτηση H48 ταχύτερη κατά τάξη μεγέθους. Επομένως οι ισχυρισμοί επιτάχυνσης δέκα φορών και άνω περιορίζονται ρητά σε περιπτώσεις κόστους φόρτωσης και ελέγχου πίνακα ή ρυθμαπόδοσης δέσμης.

\begin{table}[htbp]\centering
\resizebox{\textwidth}{!}{{\small
\begin{tabular}{lrrrr}
\hline
Case & Expected & Status & Length & Seconds \\
\hline
solved & 0 & exact & 0 & 0.0 \\
shallow\_r\_u\_f2 & 3 & exact & 3 & 8.960312 \\
deterministic\_depth\_25 & -- & exact & 18 & 56.078224 \\
superflip\_distance\_20 & 20 & exact & 20 & 158.440286 \\
\hline
\end{tabular}
}
}
\caption[Πιστοποίηση του μαντείου H48 σε έμπιστη λειτουργία με προφόρτωση]{Πιστοποίηση της διαδρομής μαντείου H48 σε έμπιστη λειτουργία με προφόρτωση (όριο 180~s): 4/4 \code{exact}, με τη γραμμή του superflip σε 158.440286~s. Πηγή: \path{results/processed/h48_oracle_certification_seed_2026_thesis_trusted_preload.json}.}
\label{tab:h48-cert-preload}
\end{table}

Καταγράφεται και τρίτη παραλλαγή της ίδιας πιστοποίησης, σε έμπιστη λειτουργία χωρίς προφόρτωση (όριο 300~s). Όπως δείχνει ο Πίνακας~\ref{tab:h48-cert-no-preload}, και εδώ οι τέσσερις γραμμές είναι \code{exact} και επαληθευμένες: η ντετερμινιστική περίπτωση ονομαστικού βάθους 25 λύνεται σε μήκος 18 σε 63.444134~s και το superflip σε μήκος 20 σε 91.629342~s. Η παραλλαγή αυτή ολοκληρώνει το δύσκολο σώμα χωρίς προφόρτωση.

\begin{table}[htbp]\centering
\resizebox{\textwidth}{!}{{\small
\begin{tabular}{lrrrr}
\hline
Case & Expected & Status & Length & Seconds \\
\hline
solved & 0 & exact & 0 & 0.0 \\
shallow\_r\_u\_f2 & 3 & exact & 3 & 2.529118 \\
deterministic\_depth\_25 & -- & exact & 18 & 63.444134 \\
superflip\_distance\_20 & 20 & exact & 20 & 91.629342 \\
\hline
\end{tabular}
}
}
\caption[Πιστοποίηση του μαντείου H48 σε έμπιστη λειτουργία χωρίς προφόρτωση]{Πιστοποίηση της διαδρομής μαντείου H48 σε έμπιστη λειτουργία χωρίς προφόρτωση (όριο 300~s): 4/4 \code{exact}, με μέγιστο χρόνο 91.629342~s για το superflip. Πηγή: \path{results/processed/h48_oracle_certification_seed_2026_thesis_trusted_no_preload.json}.}
\label{tab:h48-cert-no-preload}
\end{table}

Τέλος, η διαδρομή της παραμένουσας διεργασίας δοκιμάστηκε όχι μόνο σε ρηχές εισόδους αλλά και στο ίδιο δύσκολο σώμα πιστοποίησης. Η εκτέλεση του Πίνακα~\ref{tab:h48-resident-cert} διατηρεί μία εγγενή διεργασία H48 και έναν απεικονισμένο στη μνήμη πίνακα \code{h48h7} για όλες τις γραμμές. Όλες οι γραμμές είναι \code{exact} και επαληθευμένες: η ντετερμινιστική περίπτωση βάθους 25 λύνεται σε μήκος 18 σε 72.071228~s και το superflip σε μήκος 20 σε 77.822389~s, με συνολικό χρόνο 154.755697~s. Αυτό είναι ισχυρότερο τεκμήριο για τη βελτιστοποιημένη διαδρομή μαντείου από τις ρηχές μικρομετρήσεις, αλλά παραμένει τεκμήριο σώματος και όχι εξαντλητική απόδειξη χρόνου χειρότερης περίπτωσης.

\begin{table}[htbp]\centering
\resizebox{\textwidth}{!}{{\small
\begin{tabular}{lrrrr}
\hline
Case & Expected & Status & Length & Seconds \\
\hline
solved & 0 & exact & 0 & 0.0 \\
shallow\_r\_u\_f2 & 3 & exact & 3 & 3.49045 \\
deterministic\_depth\_25 & -- & exact & 18 & 72.071228 \\
superflip\_distance\_20 & 20 & exact & 20 & 77.822389 \\
\hline
\end{tabular}
}
}
\caption[Πιστοποίηση του δύσκολου σώματος μέσω παραμένουσας διεργασίας]{Πιστοποίηση του δύσκολου σώματος μέσω παραμένουσας εγγενούς διεργασίας με έμπιστο πίνακα \code{h48h7} (όριο 240~s): 4/4 \code{exact}, με το superflip σε 77.822389~s. Πηγή: \path{results/processed/h48_resident_certification_seed_2026_thesis_h48h7_trusted.json}.}
\label{tab:h48-resident-cert}
\end{table}

\section{Λειτουργία δέσμης και δημόσιες διεπαφές του μαντείου}
\label{sec:experiments-h48-interfaces}

Για τη μέτρηση του κόστους ανά κλήση στην ελεγχόμενη λειτουργία, 12 επαναλήψεις της ίδιας ρηχής περίπτωσης εκτελέστηκαν ως ξεχωριστές διεργασίες και ως μία δέσμη με τον πίνακα φορτωμένο μία φορά (Πίνακας~\ref{tab:h48-batch-overhead}). Όλες οι γραμμές είναι \code{exact} και επαληθευμένες: 69.740540~s συνολικά για τις ξεχωριστές διεργασίες έναντι 6.033133~s για τη δέσμη. Στην ελεγχόμενη λειτουργία κάθε ξεχωριστή διεργασία πληρώνει την πλήρη φόρτωση και τον πλήρη έλεγχο του πίνακα, οπότε η διαφορά μετρά κυρίως αυτό το επαναλαμβανόμενο κόστος· η δίκαιη ποσοτικοποίηση του καθαρού κόστους δέσμης δίνεται στην Ενότητα~\ref{sec:experiments-h48-fair}.

\begin{table}[htbp]\centering
{\small
\begin{tabular}{lrr}
\hline
Mode & Total seconds & Exact rows \\
\hline
Single process per solve & 69.74054 & 12 \\
Batch table loaded once & 6.033133 & 12 \\
Throughput speedup & 11.56 & -- \\
\hline
\end{tabular}
}

\caption[Κόστος δέσμης σε ελεγχόμενη λειτουργία χωρίς κοινή προθέρμανση]{Κόστος δέσμης σε ελεγχόμενη λειτουργία, χωρίς κοινή προθέρμανση: 12 επαναλήψεις ως ξεχωριστές διεργασίες έναντι μίας δέσμης· η διαφορά κυριαρχείται από την επαναλαμβανόμενη φόρτωση και τον επαναλαμβανόμενο έλεγχο του πίνακα και υπερεκτιμά το καθαρό όφελος της δέσμης· η δίκαιη μέτρηση δίνεται στον Πίνακα~\ref{tab:h48-batch-fair}. Πηγή: \path{results/processed/h48_batch_overhead_seed_2026_thesis.json}.}
\label{tab:h48-batch-overhead}
\end{table}

Η ίδια διαδρομή δέσμης εκτίθεται ως δημόσια εντολή \code{rubik-optimal oracle}. Η εντολή δέχεται γραμμή προς γραμμή λυμένες καταστάσεις, ακολουθίες κινήσεων ή συμβολοσειρές ψηφίδων, μετατρέπει κάθε έγκυρη κατάσταση σε άμεση είσοδο H48 και καλεί την εγγενή μηχανή δέσμης, ώστε ο πίνακας \code{h48h7} να φορτωθεί μία φορά για όλο το σύνολο εισόδων. Το αποθηκευμένο τεκμήριο της διεπαφής γραμμής εντολών (Πίνακας~\ref{tab:h48-oracle-cli}) έχει τέσσερις εισόδους, όλες \code{exact} και ανεξάρτητα επαληθευμένες. Ο συνολικός χρόνος του περιβλήματος είναι 5.917467~s, ο χρόνος δέσμης που αναφέρει η ίδια η διεπαφή είναι 5.853855~s, και οι επιμέρους αναζητήσεις των μη λυμένων γραμμών απαιτούν εκατοστά του δευτερολέπτου. Η επαναλαμβανόμενη χρήση του μαντείου επωφελείται επομένως κατά τάξη μεγέθους, επειδή το ακριβό μέρος είναι η κοινή φόρτωση και ο έλεγχος του πίνακα.

\begin{table}[htbp]\centering
{\small
\begin{tabular}{llrr}
\hline
Case & Input & Distance & Seconds \\
\hline
solved & solved & 0 & 0.000000 \\
sequence\_shallow & sequence & 3 & 0.013599 \\
facelets\_shallow & facelets & 3 & 0.005634 \\
sequence\_mixed & sequence & 5 & 0.012297 \\
\hline
\end{tabular}
}

\caption[Τεκμήριο της εντολής δέσμης του μαντείου]{Τεκμήριο της δημόσιας εντολής \code{rubik-optimal oracle} σε λειτουργία δέσμης: τέσσερις είσοδοι (λυμένη κατάσταση, ακολουθίες, ψηφίδες), όλες \code{exact}, με τις επιμέρους αναζητήσεις σε εκατοστά του δευτερολέπτου. Πηγή: τεκμήριο του \code{scripts/run_h48_oracle_cli.py} για το προφίλ \code{thesis}.}
\label{tab:h48-oracle-cli}
\end{table}

Η αντίστοιχη εκτέλεση της ίδιας διεπαφής με έμπιστο πίνακα (Πίνακας~\ref{tab:h48-oracle-cli-trusted}) διατηρεί 4/4 γραμμές \code{exact} και επαληθευμένες και μειώνει τον συνολικό χρόνο του περιβλήματος σε 2.132448~s. Η μείωση είναι μικρότερη από τη συνθετική μέτρηση ανά κλήση, επειδή η διεπαφή χρησιμοποιεί ήδη φόρτωση δέσμης· παραμένει όμως χρήσιμη πρακτική βελτίωση για τη δημόσια χρήση της εντολής.

\begin{table}[htbp]\centering
{\small
\begin{tabular}{llrr}
\hline
Case & Input & Distance & Seconds \\
\hline
solved & solved & 0 & 0.000000 \\
sequence\_shallow & sequence & 3 & 1.120128 \\
facelets\_shallow & facelets & 3 & 0.009695 \\
sequence\_mixed & sequence & 5 & 0.884141 \\
\hline
\end{tabular}
}

\caption[Η εντολή του μαντείου με έμπιστο πίνακα]{Η εντολή \code{rubik-optimal oracle} με έμπιστο πίνακα \code{h48h7}: 4/4 γραμμές \code{exact} με μειωμένο συνολικό χρόνο περιβλήματος. Πηγή: τεκμήριο του \code{scripts/run_h48_oracle_cli.py} με επιλογή έμπιστου πίνακα.}
\label{tab:h48-oracle-cli-trusted}
\end{table}

Για επαναλαμβανόμενη χρήση χωρίς αναμονή ολόκληρης δέσμης, υποστηρίζεται λειτουργία ροής: η εντολή \code{rubik-optimal oracle --stream} διατηρεί μία παραμένουσα εγγενή διεργασία H48 και τυπώνει μία γραμμή JSON μόλις ολοκληρώνεται κάθε είσοδος. Το τεκμήριο του Πίνακα~\ref{tab:h48-oracle-stream} έχει 4/4 γραμμές \code{exact} και επαληθευμένες, με λυμένη κατάσταση, ακολουθίες κινήσεων και είσοδο ψηφίδων. Η πρώτη μη λυμένη γραμμή απορροφά το κόστος ετοιμότητας της παραμένουσας διεργασίας, ενώ οι επόμενες χρησιμοποιούν τον ήδη απεικονισμένο πίνακα.

\begin{table}[htbp]\centering
{\small
\begin{tabular}{llrr}
\hline
Case & Input & Distance & Seconds \\
\hline
solved & solved & 0 & 0.000000 \\
sequence\_shallow & sequence & 3 & 1.640139 \\
facelets\_shallow & facelets & 3 & 0.012159 \\
sequence\_mixed & sequence & 5 & 2.509894 \\
\hline
\end{tabular}
}

\caption[Λειτουργία ροής του μαντείου με παραμένουσα διεργασία]{Λειτουργία ροής της εντολής \code{rubik-optimal oracle} με παραμένουσα εγγενή διεργασία και έμπιστο πίνακα: 4/4 γραμμές \code{exact}, με την πρώτη μη λυμένη γραμμή να απορροφά το κόστος ετοιμότητας. Πηγή: τεκμήριο του \code{scripts/run_h48_oracle_stream.py}.}
\label{tab:h48-oracle-stream}
\end{table}

Η ίδια υλοποίηση εκτίθεται και ως προγραμματιστική διεπαφή σε επίπεδο πακέτου, μέσω της κλάσης \code{FastOptimalOracle}. Το τεκμήριο του Πίνακα~\ref{tab:fast-oracle-api} περιλαμβάνει λυμένη κατάσταση, ακολουθία κινήσεων, είσοδο ψηφίδων και ντετερμινιστική κατάσταση ονομαστικού βάθους 10· οι 4/4 γραμμές είναι \code{exact} και επαληθευμένες, με μέγιστο χρόνο 2.664585~s. Το ίδιο αρχείο συνδέει και το τεκμήριο της πιστοποίησης του δύσκολου σώματος, ώστε η προγραμματιστική απόδειξη να μην αποκόπτεται από τις δυσκολότερες γραμμές. Και εδώ ο χρόνος εκτέλεσης παραμένει εμπειρικό τεκμήριο και όχι εξαντλητική χρονική απόδειξη για κάθε πιθανή κατάσταση.

\begin{table}[htbp]\centering
{\small
\begin{tabular}{llrrr}
\hline
Case & Input & Expected & Length & Seconds \\
\hline
solved & cube\_state & 0 & 0 & 0.0 \\
sequence\_shallow & sequence & 3 & 3 & 2.393554 \\
facelets\_shallow & facelets & 3 & 3 & 0.020547 \\
deterministic\_depth\_10 & deterministic\_scramble & -- & 10 & 2.664585 \\
\hline
\end{tabular}
}

\caption[Τεκμήριο της προγραμματιστικής διεπαφής του μαντείου]{Τεκμήριο της προγραμματιστικής διεπαφής \code{FastOptimalOracle} με έμπιστο πίνακα \code{h48h7}: 4/4 γραμμές \code{exact} με μέγιστο χρόνο 2.664585~s. Πηγή: τεκμήριο του \code{scripts/run_fast_optimal_oracle_api.py}.}
\label{tab:fast-oracle-api}
\end{table}

\section{Συμβόλαιο ισχυρισμών}
\label{sec:experiments-h48-contract}

Για να μην παραμείνει το όριο των ισχυρισμών μόνο ως αφηγηματικό σχόλιο, παράγεται χωριστό τεκμήριο συμβολαίου που ελέγχει αυτόματα τις κρίσιμες ιδιότητες της διαδρομής μαντείου. Οι ουσιώδεις έλεγχοι είναι πέντε. Πρώτον, η τεκμηρίωση και οι επικεφαλίδες του ενσωματωμένου nissy-core επιβεβαιώνουν ότι η διαδρομή H48 είναι βέλτιστος επιλυτής στη μετρική μισών στροφών και ότι το ψευδώνυμο \code{optimal} δείχνει στο \code{h48h7}. Δεύτερον, το περίβλημα C καλεί τη \code{nissy_solve} με σημαία κανονικής λειτουργίας και μέγιστο βάθος 20. Τρίτον, το περίβλημα Python και η κλάση \code{FastOptimalOracle} εκτελούν άμεση μετατροπή κατάστασης κυβιδίων (\eng{cubies}), έλεγχο φυσικής εγκυρότητας, χρήση παραμένουσας μηχανής \code{h48h7} με μεταδεδομένα παραγωγής και ανεξάρτητη επαλήθευση κάθε λύσης. Τέταρτον, το συμβόλαιο συνδέει τα αποθηκευμένα εμπειρικά τεκμήρια του παραρτήματος και του Κεφαλαίου~\ref{chap:experiments}. Πέμπτον, το πεδίο της καθολικής απόδειξης ταχείας εκτέλεσης για κάθε κατάσταση παραμένει σκόπιμα αρνητικό: η διεπαφή δέχεται κάθε φυσικά έγκυρη κατάσταση 3×3×3 και τηρεί σύμβαση αποκλειστικά ακριβών αποτελεσμάτων για όσες γραμμές ολοκληρώνονται, χωρίς να μετατρέπει τους χρόνους του σώματος σε θεώρημα χρόνου χειρότερης περίπτωσης. Ο πλήρης πίνακας των ελέγχων του συμβολαίου παρατίθεται στο Παράρτημα~\ref{app:repro}.

\section{Χαρτοφυλάκιο ακριβών επιλυτών}
\label{sec:experiments-portfolio}

Η πρακτική συνέπεια της ύπαρξης δύο ακριβών μηχανών είναι ότι δεν χρειάζεται να επιλεγεί ένας μοναδικός επιλυτής για όλες τις περιπτώσεις. Η κλάση \code{PortfolioOptimalOracle} επαναχρησιμοποιεί επανεπικυρωμένα ακριβή πιστοποιητικά για ήδη αποδεδειγμένες καταστάσεις, δοκιμάζει πρώτα τη διαδρομή \code{nissy-optimal} με περιορισμένο χρόνο και, αν αυτή δεν επιστρέψει ακριβές επαληθευμένο αποτέλεσμα, μεταβιβάζει την ίδια κατάσταση στο παραμένον μαντείο H48. Η διαδρομή αυτή δεν αλλάζει τον κανόνα ακρίβειας: η κατάσταση \code{exact} αποδίδεται μόνο όταν η επιλεγμένη μηχανή επιστρέψει λύση που επαληθεύεται ανεξάρτητα, ή όταν ένα αποθηκευμένο ακριβές και επαληθευμένο πιστοποιητικό επανελεγχθεί πάνω στην ίδια κατάσταση.

Σε εκτέλεση χαμηλής προτεραιότητας υπό φορτίο γενικής χρήσης, το σώμα με προτεραιότητα στο Nissy χρησιμοποιεί μία διεργασία δέσμης Nissy, έχει 5/5 γραμμές \code{exact} και μέγιστο χρόνο δέσμης 10.097333~s (Πίνακας~\ref{tab:portfolio-nissy-first}). Η στοχευμένη γραμμή εφεδρείας του superflip λύνει το superflip απόστασης 20 μέσω παραμένουσας διεργασίας H48 σε 135.679618~s με 4 νήματα H48 (Πίνακας~\ref{tab:portfolio-superflip-fallback}). Μετά την αποθήκευση του αντίστοιχου ακριβούς πιστοποιητικού, η ίδια κατάσταση επιστρέφεται από την κρυφή μνήμη πιστοποιητικών σε 0.012885~s χωρίς νέα ακριβή αναζήτηση (Πίνακας~\ref{tab:portfolio-superflip-cache}). Το αποτέλεσμα είναι τεκμήριο μηχανικής σύνθεσης χαρτοφυλακίου, όχι απόδειξη ότι κάθε πιθανή κατάσταση έχει πρακτικό χρόνο χειρότερης περίπτωσης.

\begin{table}[htbp]\centering
{\small
\begin{tabular}{lrrrr}
\hline
Case & Source depth & Status & Length & Seconds \\
\hline
solved & 0 & exact & 0 & 0.0 \\
shallow\_sequence & 3 & exact & 3 & 10.097333 \\
random\_1\_10 & 10 & exact & 9 & 10.097333 \\
random\_2\_15 & 15 & exact & 14 & 10.097333 \\
random\_3\_20 & 20 & exact & 17 & 10.097333 \\
\hline
\end{tabular}
}

\caption[Σώμα του χαρτοφυλακίου με προτεραιότητα στο Nissy]{Σώμα του χαρτοφυλακίου ακριβών επιλυτών με προτεραιότητα στο Nissy, σε εκτέλεση χαμηλής προτεραιότητας: 5/5 γραμμές \code{exact} μέσω μίας διεργασίας δέσμης. Πηγή: τεκμήριο \code{lowload} του \code{PortfolioOptimalOracle}.}
\label{tab:portfolio-nissy-first}
\end{table}

\begin{table}[htbp]\centering
{\small
\begin{tabular}{lrrrr}
\hline
Case & Source depth & Status & Length & Seconds \\
\hline
superflip\_distance\_20 & 0 & exact & 20 & 135.679618 \\
\hline
\end{tabular}
}

\caption[Γραμμή εφεδρείας του χαρτοφυλακίου για το superflip]{Γραμμή εφεδρείας του χαρτοφυλακίου για το superflip: ακριβής λύση μήκους 20 μέσω παραμένουσας διεργασίας H48 σε 135.679618~s. Πηγή: τεκμήριο \code{lowload} του \code{PortfolioOptimalOracle}.}
\label{tab:portfolio-superflip-fallback}
\end{table}

\begin{table}[htbp]\centering
{\small
\begin{tabular}{lrrrr}
\hline
Case & Source depth & Status & Length & Seconds \\
\hline
superflip\_distance\_20 & 0 & exact & 20 & 0.012885 \\
\hline
\end{tabular}
}

\caption[Επιστροφή του superflip από την κρυφή μνήμη πιστοποιητικών]{Επιστροφή του superflip από την κρυφή μνήμη πιστοποιητικών του χαρτοφυλακίου σε 0.012885~s, μετά την επανεπικύρωση του αποθηκευμένου ακριβούς πιστοποιητικού. Πηγή: τεκμήριο \code{lowload} του \code{PortfolioOptimalOracle}.}
\label{tab:portfolio-superflip-cache}
\end{table}

\section{Η δοκιμή παραγωγής του πίνακα h48h8}
\label{sec:h48eng-h48h8}

Δοκιμάστηκε η παραγωγή του μεγαλύτερου πίνακα \code{h48h8} στο ίδιο μηχάνημα. Η προσπάθεια διακόπηκε μετά από εκτέλεση άνω των δύο ωρών χωρίς να παραχθεί το αρχείο \code{h48h8.bin}, με δείγματα εκτέλεσης να δείχνουν ότι ο χρόνος καταναλωνόταν μέσα στον εγγενή γεννήτορα H48 και όχι στο στρώμα Python. Επειδή δεν υπάρχει πλήρες αρχείο, άθροισμα ελέγχου ή μεταδεδομένα, η εργασία δεν χρησιμοποιεί το \code{h48h8} ως αποτέλεσμα και δεν αποδίδει σε αυτό κανέναν ισχυρισμό απόδοσης.
